\section{Local complete intersections}
\label{section-lci}

\noindent
The property of being a local complete intersection is an
intrinsic property of a Noetherian local ring.
This will be discussed in
Divided Power Algebra, Section \ref{dpa-section-lci}.
However, for the moment we just define this property for
finite type algebras over a field.

\begin{definition}
\label{definition-lci-field}
Let $k$ be a field.
Let $S$ be a finite type $k$-algebra.
\begin{enumerate}
\item We say that $S$ is a {\it global complete intersection over $k$}
if there exists a presentation $S = k[x_1, \ldots, x_n]/(f_1, \ldots, f_c)$
such that $\dim(S) = n - c$.
\item We say that $S$ is a {\it local complete intersection over $k$}
if there exists a covering $\Spec(S) = \bigcup D(g_i)$ such
that each of the rings $S_{g_i}$ is a global complete intersection
over $k$.
\end{enumerate}
We will also use the convention that the zero ring is a global
complete intersection over $k$.
\end{definition}

\noindent
Suppose $S$ is a global complete intersection
$S = k[x_1, \ldots, x_n]/(f_1, \ldots, f_c)$ as in
Definition \ref{definition-lci-field}.
For a maximal ideal $\mathfrak m \subset k[x_1, \ldots, x_n]$
we have $\dim(k[x_1, \ldots, x_n]_\mathfrak m) = n$
(Lemma \ref{lemma-dim-affine-space}).
If $(f_1, \ldots, f_c) \subset \mathfrak m$, then we
conclude that $\dim(S_\mathfrak m) \geq n - c$ by
Lemma \ref{lemma-one-equation}. Since $\dim(S) = n - c$
by Definition \ref{definition-lci-field} we conclude
that $\dim(S_\mathfrak m) = n - c$ for all maximal ideals of $S$
and that $\Spec(S)$ is equidimensional
(Topology, Definition \ref{topology-definition-equidimensional})
of dimension $n - c$, see
Lemma \ref{lemma-dimension-at-a-point-finite-type-over-field}.
We will often use this without further mention.

\begin{lemma}
\label{lemma-localize-lci}
Let $k$ be a field.
Let $S$ be a finite type $k$-algebra.
Let $g \in S$.
\begin{enumerate}
\item If $S$ is a global complete intersection so is $S_g$.
\item If $S$ is a local complete intersection so is $S_g$.
\end{enumerate}
\end{lemma}

\begin{proof}
The second statement follows immediately from the first.
Proof of the first statement. If $S_g$ is the zero ring,
then it is true. Assume $S_g$ is nonzero.
Write $S = k[x_1, \ldots, x_n]/(f_1, \ldots, f_c)$
with $n - c = \dim(S)$ as in Definition \ref{definition-lci-field}.
By the remarks following the definition $\dim(S_g) = n - c$. Let
$g' \in k[x_1, \ldots, x_n]$ be an element whose residue class
corresponds to $g$. Then
$S_g =  k[x_1, \ldots, x_n, x_{n + 1}]/(f_1, \ldots, f_c, x_{n + 1}g' - 1)$
as desired.
\end{proof}

\begin{lemma}
\label{lemma-lci-CM}
Let $k$ be a field. Let $S$ be a finite type $k$-algebra.
If $S$ is a local complete intersection, then
$S$ is a Cohen-Macaulay ring.
\end{lemma}

\begin{proof}
Choose a maximal prime $\mathfrak m$ of $S$.
We have to show that $S_\mathfrak m$ is Cohen-Macaulay.
By assumption we may assume $S = k[x_1, \ldots, x_n]/(f_1, \ldots, f_c)$
with $\dim(S) = n - c$. Let $\mathfrak m' \subset k[x_1, \ldots, x_n]$
be the maximal ideal corresponding to $\mathfrak m$.
According to Proposition \ref{proposition-finite-gl-dim-polynomial-ring}
the local ring
$k[x_1, \ldots, x_n]_{\mathfrak m'}$ is regular local of
dimension $n$. In particular it is Cohen-Macaulay by
Lemma \ref{lemma-regular-ring-CM}.
By Lemma \ref{lemma-one-equation} applied $c$ times the local ring
$S_{\mathfrak m} = k[x_1, \ldots, x_n]_{\mathfrak m'}/(f_1, \ldots, f_c)$
has dimension $\geq n - c$. By assumption $\dim(S_{\mathfrak m}) \leq n - c$.
Thus we get equality. This implies that $f_1, \ldots, f_c$ is a regular
sequence in $k[x_1, \ldots, x_n]_{\mathfrak m'}$ and that
$S_{\mathfrak m}$ is Cohen-Macaulay, see Proposition
\ref{proposition-CM-module}.
\end{proof}

\noindent
The following is the technical key to the rest of the material in this
section. An important feature of this lemma is that we may choose any
presentation for the ring $S$, but that condition (1) does not depend
on this choice.

\begin{lemma}
\label{lemma-lci}
Let $k$ be a field.
Let $S$ be a finite type $k$-algebra.
Let $\mathfrak q$ be a prime of $S$.
Choose any presentation $S = k[x_1, \ldots, x_n]/I$.
Let $\mathfrak q'$ be the prime of $k[x_1, \ldots, x_n]$ corresponding
to $\mathfrak q$. Set
$c = \text{height}(\mathfrak q') - \text{height}(\mathfrak q)$,
in other words $\dim_{\mathfrak q}(S) = n - c$
(see Lemma \ref{lemma-codimension}). The following are equivalent
\begin{enumerate}
\item There exists a $g \in S$, $g \not \in \mathfrak q$
such that $S_g$ is a global complete intersection over $k$.
\item The ideal $I_{\mathfrak q'} \subset k[x_1, \ldots, x_n]_{\mathfrak q'}$
can be generated by $c$ elements.
\item The conormal module $(I/I^2)_{\mathfrak q}$ can be generated by
$c$ elements over $S_{\mathfrak q}$.
\item The conormal module $(I/I^2)_{\mathfrak q}$ is a free
$S_{\mathfrak q}$-module of rank $c$.
\item The ideal $I_{\mathfrak q'}$ can be generated by a regular sequence
in the regular local ring $k[x_1, \ldots, x_n]_{\mathfrak q'}$.
\end{enumerate}
In this case any $c$ elements of $I_{\mathfrak q'}$
which generate $I_{\mathfrak q'}/\mathfrak q'I_{\mathfrak q'}$
form a regular sequence in the local
ring $k[x_1, \ldots, x_n]_{\mathfrak q'}$.
\end{lemma}

\begin{proof}
Set $R = k[x_1, \ldots, x_n]_{\mathfrak q'}$. This is a
Cohen-Macaulay local
ring of dimension $\text{height}(\mathfrak q')$, see for example
Lemma \ref{lemma-lci-CM}. Moreover,
$\overline{R} = R/IR = R/I_{\mathfrak q'} = S_{\mathfrak q}$
is a quotient of dimension $\text{height}(\mathfrak q)$.
Let $f_1, \ldots, f_c \in I_{\mathfrak q'}$ be elements
which generate $(I/I^2)_{\mathfrak q}$. By Lemma \ref{lemma-NAK}
we see that $f_1, \ldots, f_c$ generate $I_{\mathfrak q'}$.
Since the dimensions work out, we conclude
by Proposition \ref{proposition-CM-module} that
$f_1, \ldots, f_c$ is a regular sequence in $R$.
By Lemma \ref{lemma-regular-quasi-regular} we see that
$(I/I^2)_{\mathfrak q}$ is free.
These arguments show that (2), (3), (4) are equivalent and
that they imply the last statement of the lemma, and therefore
they imply (5).

\medskip\noindent
If (5) holds, say $I_{\mathfrak q'}$ is generated by a regular
sequence of length $e$, then
$\text{height}(\mathfrak q) = \dim(S_{\mathfrak q}) =
\dim(k[x_1, \ldots, x_n]_{\mathfrak q'}) - e =
\text{height}(\mathfrak q') - e$ by dimension theory,
see Section \ref{section-dimension}. We conclude that $e = c$.
Thus (5) implies (2).

\medskip\noindent
We continue with the notation introduced in the first paragraph.
For each $f_i$ we may find $d_i \in k[x_1, \ldots, x_n]$,
$d_i \not \in \mathfrak q'$ such that
$f_i' = d_i f_i \in k[x_1, \ldots, x_n]$.
Then it is still true that $I_{\mathfrak q'} = (f_1', \ldots, f_c')R$.
Hence there exists a $g' \in k[x_1, \ldots, x_n]$, $g' \not \in \mathfrak q'$
such that $I_{g'} = (f_1', \ldots, f_c')$.
Moreover, pick $g'' \in k[x_1, \ldots, x_n]$, $g'' \not \in \mathfrak q'$
such that $\dim(S_{g''}) = \dim_{\mathfrak q} \Spec(S)$.
By Lemma \ref{lemma-codimension} this dimension is equal to $n - c$.
Finally, set $g$ equal to the image of $g'g''$ in $S$.
Then we see that
$$
S_g \cong k[x_1, \ldots, x_n, x_{n + 1}]
/
(f_1', \ldots, f_c', x_{n + 1}g'g'' - 1)
$$
and by our choice of $g''$ this ring has dimension $n - c$.
Therefore it is a global complete intersection.
Thus each of (2), (3), and (4) implies (1).

\medskip\noindent
Assume (1). Let $S_g \cong k[y_1, \ldots, y_m]/(f_1, \ldots, f_t)$
be a presentation of $S_g$ as a global complete intersection.
Write $J = (f_1, \ldots, f_t)$. Let $\mathfrak q'' \subset k[y_1, \ldots, y_m]$
be the prime corresponding to $\mathfrak qS_g$. Note that
$t = m - \dim(S_g) =
\text{height}(\mathfrak q'') - \text{height}(\mathfrak q)$,
see Lemma \ref{lemma-codimension} for the last equality.
As seen in the proof of Lemma \ref{lemma-lci-CM} (and also above) the elements
$f_1, \ldots, f_t$ form a regular sequence in the local ring
$k[y_1, \ldots, y_m]_{\mathfrak q''}$.
By Lemma \ref{lemma-regular-quasi-regular} we see that
$(J/J^2)_{\mathfrak q}$ is free of rank $t$.
By Lemma \ref{lemma-conormal-module-localize} we have
$$
J/J^2 \oplus S_g^n \cong (I/I^2)_g \oplus S_g^m
$$
Thus $(I/I^2)_{\mathfrak q}$ is free of rank
$t + n - m = m - \dim(S_g) + n - m = n - \dim(S_g) =
\text{height}(\mathfrak q') - \text{height}(\mathfrak q) = c$.
Thus we obtain (4).
\end{proof}

\noindent
The result of Lemma \ref{lemma-lci} suggests the following definition.

\begin{definition}
\label{definition-lci-local-ring}
Let $k$ be a field. Let $S$ be a local $k$-algebra essentially of finite type
over $k$. We say $S$ is a {\it complete intersection (over $k$)}
if there exists a local $k$-algebra $R$ and elements
$f_1, \ldots, f_c \in \mathfrak m_R$ such that
\begin{enumerate}
\item $R$ is essentially of finite type over $k$,
\item $R$ is a regular local ring,
\item $f_1, \ldots, f_c$ form a regular sequence in $R$, and
\item $S \cong R/(f_1, \ldots, f_c)$ as $k$-algebras.
\end{enumerate}
\end{definition}

\noindent
By the Cohen structure theorem (see
Theorem \ref{theorem-cohen-structure-theorem}) any complete
Noetherian local ring may be written as the quotient of some regular complete
local ring. Hence we may use the definition above to define the notion of
a complete intersection ring for any complete Noetherian local ring.
We will discuss this in
Divided Power Algebra, Section \ref{dpa-section-lci}.
In the meantime the following lemma shows that such a definition makes sense.

\begin{lemma}
\label{lemma-ci-well-defined}
Let $A \to B \to C$ be surjective local ring homomorphisms.
Assume $A$ and $B$ are regular local rings. The following are equivalent
\begin{enumerate}
\item $\Ker(A \to C)$ is generated by a regular sequence,
\item $\Ker(A \to C)$ is generated by $\dim(A) - \dim(C)$ elements,
\item $\Ker(B \to C)$ is generated by a regular sequence, and
\item $\Ker(B \to C)$ is generated by $\dim(B) - \dim(C)$ elements.
\end{enumerate}
\end{lemma}

\begin{proof}
A regular local ring is Cohen-Macaulay, see Lemma \ref{lemma-regular-ring-CM}.
Hence the equivalences (1) $\Leftrightarrow$ (2) and
(3) $\Leftrightarrow$ (4), see Proposition \ref{proposition-CM-module}.
By Lemma \ref{lemma-regular-quotient-regular}
the ideal $\Ker(A \to B)$ can be generated
by $\dim(A) - \dim(B)$ elements.
Hence we see that (4) implies (2).

\medskip\noindent
It remains to show that (1) implies (4). We do this by induction on
$\dim(A) - \dim(B)$. The case $\dim(A) - \dim(B) = 0$ is trivial.
Assume $\dim(A) > \dim (B)$.
Write $I = \Ker(A \to C)$ and $J = \Ker(A \to B)$.
Note that $J \subset I$. Our assumption is that the minimal number
of generators of $I$ is $\dim(A) - \dim(C)$.
Let $\mathfrak m \subset A$ be the maximal
ideal. Consider the maps
$$
J/ \mathfrak m J \to  I / \mathfrak m I \to \mathfrak m /\mathfrak m^2
$$
By Lemma \ref{lemma-regular-quotient-regular} and its proof the
composition is injective. Take any element $x \in J$ which is
not zero in $J /\mathfrak mJ$. By the above and Nakayama's lemma
$x$ is an element of a minimal set of generators of $I$.
Hence we may replace $A$ by $A/xA$ and $I$ by $I/xA$ which
decreases both $\dim(A)$ and the minimal number of generators of $I$
by $1$. Thus we win.
\end{proof}

\begin{lemma}
\label{lemma-lci-local}
Let $k$ be a field. Let $S$ be a local $k$-algebra essentially of finite
type over $k$. The following are equivalent:
\begin{enumerate}
\item $S$ is a complete intersection over $k$,
\item for any surjection $R \to S$ with $R$ a regular local ring
essentially of finite presentation over $k$ the ideal
$\Ker(R \to S)$ can be generated by a regular sequence,
\item for some surjection $R \to S$ with $R$ a regular local ring
essentially of finite presentation over $k$ the ideal
$\Ker(R \to S)$ can be generated by
$\dim(R) - \dim(S)$ elements,
\item there exists a global complete intersection
$A$ over $k$ and a prime $\mathfrak a$ of $A$ such
that $S \cong A_{\mathfrak a}$, and
\item there exists a local complete intersection
$A$ over $k$ and a prime $\mathfrak a$ of $A$ such
that $S \cong A_{\mathfrak a}$.
\end{enumerate}
\end{lemma}

\begin{proof}
It is clear that (2) implies (1) and (1) implies (3).
It is also clear that (4) implies (5). Let us show that (3) implies
(4). Thus we assume there exists a surjection
$R \to S$ with $R$ a regular local ring
essentially of finite presentation over $k$ such that the ideal
$\Ker(R \to S)$ can be generated by $\dim(R) - \dim(S)$ elements.
We may write $R = (k[x_1, \ldots, x_n]/J)_{\mathfrak q}$
for some $J \subset k[x_1, \ldots, x_n]$ and
some prime $\mathfrak q \subset k[x_1, \ldots, x_n]$ with
$J \subset \mathfrak q$. Let $I \subset k[x_1, \ldots, x_n]$
be the kernel of the map $k[x_1, \ldots, x_n] \to S$ so that
$S \cong (k[x_1, \ldots, x_n]/I)_{\mathfrak q}$.
By assumption $(I/J)_{\mathfrak q}$ is generated by
$\dim(R) - \dim(S)$ elements. We conclude that
$I_{\mathfrak q}$ can be generated by
$\dim(k[x_1, \ldots, x_n]_{\mathfrak q}) - \dim(S)$ elements
by Lemma \ref{lemma-ci-well-defined}.
From Lemma \ref{lemma-lci} we see that for some
$g \in k[x_1, \ldots, x_n]$, $g \not \in \mathfrak q$
the algebra $(k[x_1, \ldots, x_n]/I)_g$ is a global
complete intersection and $S$ is isomorphic to
a local ring of it.

\medskip\noindent
To finish the proof of the lemma we have to show that (5) implies (2).
Assume (5) and let $\pi : R \to S$ be a surjection with $R$ a regular local
$k$-algebra essentially of finite type over $k$.
By assumption we have $S = A_{\mathfrak a}$ for some local
complete intersection $A$ over $k$.
Choose a presentation $R = (k[y_1, \ldots, y_m]/J)_{\mathfrak q}$
with $J \subset \mathfrak q \subset k[y_1, \ldots, y_m]$.
We may and do assume that $J$ is the kernel of the map
$k[y_1, \ldots, y_m] \to R$. Let $I \subset k[y_1, \ldots, y_m]$
be the kernel of the map $k[y_1, \ldots, y_m] \to S = A_{\mathfrak a}$.
Then $J \subset I$ and $(I/J)_{\mathfrak q}$ is the kernel of
the surjection $\pi : R \to S$. So
$S = (k[y_1, \ldots, y_m]/I)_{\mathfrak q}$.

\medskip\noindent
By Lemma \ref{lemma-isomorphic-local-rings} we see that there exist
$g \in A$, $g \not \in \mathfrak a$ and
$g' \in k[y_1, \ldots, y_m]$, $g' \not \in \mathfrak q$
such that $A_g \cong (k[y_1, \ldots, y_m]/I)_{g'}$.
After replacing $A$ by $A_g$ and $k[y_1, \ldots, y_m]$ by
$k[y_1, \ldots, y_{m + 1}]$ we may assume that
$A \cong k[y_1, \ldots, y_m]/I$. Consider the surjective
maps of local rings
$$
k[y_1, \ldots, y_m]_{\mathfrak q} \to R \to S.
$$
We have to show that the kernel of $R \to S$ is generated by
a regular sequence. By Lemma \ref{lemma-lci} we know that
$k[y_1, \ldots, y_m]_{\mathfrak q} \to A_{\mathfrak a} = S$
has this property (as $A$ is a local complete intersection over $k$).
We win by Lemma \ref{lemma-ci-well-defined}.
\end{proof}

\begin{lemma}
\label{lemma-lci-at-prime}
Let $k$ be a field. Let $S$ be a finite type $k$-algebra.
Let $\mathfrak q$ be a prime of $S$. The following are
equivalent:
\begin{enumerate}
\item The local ring $S_{\mathfrak q}$ is a complete intersection
ring (Definition \ref{definition-lci-local-ring}).
\item There exists a $g \in S$, $g \not \in \mathfrak q$
such that $S_g$ is a local complete intersection over $k$.
\item There exists a $g \in S$, $g \not \in \mathfrak q$
such that $S_g$ is a global complete intersection over $k$.
\item For any presentation $S = k[x_1, \ldots, x_n]/I$ with
$\mathfrak q' \subset k[x_1, \ldots, x_n]$ corresponding to $\mathfrak q$
any of the equivalent conditions (1) -- (5) of Lemma \ref{lemma-lci} hold.
\end{enumerate}
\end{lemma}

\begin{proof}
This is a combination of Lemmas \ref{lemma-lci} and \ref{lemma-lci-local}
and the definitions.
\end{proof}

\begin{lemma}
\label{lemma-lci-global}
Let $k$ be a field. Let $S$ be a finite type $k$-algebra.
The following are equivalent:
\begin{enumerate}
\item The ring $S$ is a local complete intersection over $k$.
\item All local rings of $S$ are complete intersection rings over $k$.
\item All localizations of $S$
at maximal ideals are complete intersection rings over $k$.
\end{enumerate}
\end{lemma}

\begin{proof}
This follows from Lemma \ref{lemma-lci-at-prime},
the fact that $\Spec(S)$ is quasi-compact and the definitions.
\end{proof}

\noindent
The following lemma says that being a complete intersection is
preserved under change of base field (in a strong sense).

\begin{lemma}
\label{lemma-lci-field-change-local}
Let $K/k$ be a field extension.
Let $S$ be a finite type algebra over $k$.
Let $\mathfrak q_K$ be a prime of $S_K = K \otimes_k S$
and let $\mathfrak q$ be the corresponding prime of $S$.
Then $S_{\mathfrak q}$ is a complete intersection
over $k$ (Definition \ref{definition-lci-local-ring})
if and only if $(S_K)_{\mathfrak q_K}$ is a complete
intersection over $K$.
\end{lemma}

\begin{proof}
Choose a presentation $S = k[x_1, \ldots, x_n]/I$.
This gives a presentation
$S_K = K[x_1, \ldots, x_n]/I_K$ where $I_K = K \otimes_k I$.
Let $\mathfrak q_K' \subset K[x_1, \ldots, x_n]$,
resp.\ $\mathfrak q' \subset k[x_1, \ldots, x_n]$ be
the corresponding prime. We will show that the equivalent conditions
of Lemma \ref{lemma-lci}
hold for the pair $(S = k[x_1, \ldots, x_n]/I, \mathfrak q)$
if and only if they hold for the pair
$(S_K = K[x_1, \ldots, x_n]/I_K, \mathfrak q_K)$.
The lemma will follow from this (see Lemma \ref{lemma-lci-at-prime}).

\medskip\noindent
By Lemma \ref{lemma-dimension-at-a-point-preserved-field-extension} we have
$\dim_{\mathfrak q} S = \dim_{\mathfrak q_K} S_K$.
Hence the integer $c$ occurring in Lemma \ref{lemma-lci}
is the same for the pair $(S = k[x_1, \ldots, x_n]/I, \mathfrak q)$
as for the pair $(S_K = K[x_1, \ldots, x_n]/I_K, \mathfrak q_K)$.
On the other hand we have
\begin{eqnarray*}
I \otimes_{k[x_1, \ldots, x_n]} \kappa(\mathfrak q')
\otimes_{\kappa(\mathfrak q')} \kappa(\mathfrak q_K')
& = &
I \otimes_{k[x_1, \ldots, x_n]} \kappa(\mathfrak q_K') \\
& = &
I \otimes_{k[x_1, \ldots, x_n]} K[x_1, \ldots, x_n]
\otimes_{K[x_1, \ldots, x_n]} \kappa(\mathfrak q_K') \\
& = &
(K \otimes_k I) \otimes_{K[x_1, \ldots, x_n]} \kappa(\mathfrak q_K') \\
& = &
I_K \otimes_{K[x_1, \ldots, x_n]} \kappa(\mathfrak q'_K).
\end{eqnarray*}
Therefore,
$\dim_{\kappa(\mathfrak q')}
I \otimes_{k[x_1, \ldots, x_n]} \kappa(\mathfrak q')
=
\dim_{\kappa(\mathfrak q'_K)}
I_K \otimes_{K[x_1, \ldots, x_n]} \kappa(\mathfrak q_K')$.
Thus it follows from
Nakayama's Lemma \ref{lemma-NAK} that the minimal number
of generators of $I_{\mathfrak q'}$ is the same as the minimal
number of generators of $(I_K)_{\mathfrak q'_K}$.
Thus the lemma follows from characterization (2) of Lemma \ref{lemma-lci}.
\end{proof}

\begin{lemma}
\label{lemma-lci-field-change}
Let $k \to K$ be a field extension.
Let $S$ be a finite type $k$-algebra.
Then $S$ is a local complete intersection over $k$ if and
only if $S \otimes_k K$ is a local complete intersection over $K$.
\end{lemma}

\begin{proof}
This follows from a combination of Lemmas
\ref{lemma-lci-global} and \ref{lemma-lci-field-change-local}.
But we also give a different
proof here (based on the same principles).

\medskip\noindent
Set $S' = S \otimes_k K$. Let $\alpha : k[x_1, \ldots, x_n] \to S$ be a
presentation with kernel $I$. Let $\alpha' : K[x_1, \ldots, x_n] \to S'$
be the induced presentation with kernel $I'$.

\medskip\noindent
Suppose that $S$ is a local complete intersection.
Pick a prime $\mathfrak q \subset S'$. Denote
$\mathfrak q'$ the corresponding prime of $K[x_1, \ldots, x_n]$,
$\mathfrak p$ the corresponding prime of $S$, and
$\mathfrak p'$ the corresponding prime of $k[x_1, \ldots, x_n]$.
Consider the following diagram of Noetherian local rings
$$
\xymatrix{
S'_{\mathfrak q} &  K[x_1, \ldots, x_n]_{\mathfrak q'} \ar[l] \\
S_{\mathfrak p}\ar[u] &  k[x_1, \ldots, x_n]_{\mathfrak p'} \ar[u] \ar[l]
}
$$
By Lemma \ref{lemma-lci} we know that $S_{\mathfrak p}$
is cut out by some regular sequence $f_1, \ldots, f_c$ in
$k[x_1, \ldots, x_n]_{\mathfrak p'}$. Since the right vertical
arrow is flat we see that the images of $f_1, \ldots, f_c$
form a regular sequence in $K[x_1, \ldots, x_n]_{\mathfrak q'}$.
Because tensoring with $K$ over $k$ is an exact functor we have
$S'_{\mathfrak q} = K[x_1, \ldots, x_n]_{\mathfrak q'}/(f_1, \ldots, f_c)$.
Hence by Lemma \ref{lemma-lci} again we see that $S'$ is a local
complete intersection in a neighbourhood of $\mathfrak q$. Since
$\mathfrak q$ was arbitrary we see that $S'$ is a local complete
intersection over $K$.

\medskip\noindent
Suppose that $S'$ is a local complete intersection.
Pick a maximal ideal $\mathfrak m$ of $S$. Let $\mathfrak m'$
denote the corresponding maximal ideal of $k[x_1, \ldots, x_n]$.
Denote $\kappa = \kappa(\mathfrak m)$ the residue field.
By Remark \ref{remark-fundamental-diagram} the primes of
$S'$ lying over $\mathfrak m$ correspond to primes
in $K \otimes_k \kappa$. By the Hilbert-Nullstellensatz
Theorem \ref{theorem-nullstellensatz} we have $[\kappa : k] < \infty$.
Hence $K \otimes_k \kappa$ is finite nonzero over $K$.
Hence $K \otimes_k \kappa$ has a finite number $> 0$ of primes
which are all maximal, each of which has a residue field
finite over $K$ (see Section \ref{section-artinian}).
Hence there are finitely many $> 0$ prime ideals
$\mathfrak n \subset S'$ lying over $\mathfrak m$,
each of which is maximal and  has a residue field
which is finite over $K$. Pick one, say $\mathfrak n \subset S'$,
and let $\mathfrak n' \subset K[x_1, \ldots, x_n]$ denote the corresponding
prime ideal of $K[x_1, \ldots, x_n]$.
Note that since $V(\mathfrak mS')$ is finite, we see that
$\mathfrak n$ is an isolated closed point of it, and we
deduce that $\mathfrak mS'_{\mathfrak n}$ is an ideal of definition
of $S'_{\mathfrak n}$. This implies that
$\dim(S_{\mathfrak m}) = \dim(S'_{\mathfrak n})$ for example by
Lemma \ref{lemma-dimension-base-fibre-equals-total}.
(This can also be seen using
Lemma \ref{lemma-dimension-at-a-point-preserved-field-extension}.)
Consider the corresponding diagram of Noetherian local rings
$$
\xymatrix{
S'_{\mathfrak n} &  K[x_1, \ldots, x_n]_{\mathfrak n'} \ar[l] \\
S_{\mathfrak m}\ar[u] &  k[x_1, \ldots, x_n]_{\mathfrak m'} \ar[u] \ar[l]
}
$$
According to Lemma \ref{lemma-change-base-NL} we have
$\NL(\alpha) \otimes_S S' = \NL(\alpha')$, in particular
$I'/(I')^2 = I/I^2 \otimes_S S'$. Thus
$(I/I^2)_{\mathfrak m} \otimes_{S_{\mathfrak m}} \kappa$
and
$(I'/(I')^2)_{\mathfrak n} \otimes_{S'_{\mathfrak n}} \kappa(\mathfrak n)$
have the same dimension. Since $(I'/(I')^2)_{\mathfrak n}$
is free of rank $n - \dim S'_{\mathfrak n}$ we deduce that
$(I/I^2)_{\mathfrak m}$ can be generated by
$n - \dim S'_{\mathfrak n} = n - \dim S_{\mathfrak m}$ elements.
By Lemma \ref{lemma-lci} we see that $S$ is a local
complete intersection in a neighbourhood of $\mathfrak m$.
Since $\mathfrak m$ was any maximal ideal we conclude that
$S$ is a local complete intersection.
\end{proof}

\noindent
We end with a lemma which we will later use to prove that
given ring maps $T \to A \to B$ where $B$ is syntomic over $T$,
and $B$ is syntomic over $A$, then $A$ is syntomic over $T$.

\begin{lemma}
\label{lemma-lci-permanence-initial}
Let
$$
\xymatrix{
B & S \ar[l] \\
A \ar[u] & R \ar[l] \ar[u]
}
$$
be a commutative square of local rings. Assume
\begin{enumerate}
\item $R$ and $\overline{S} = S/\mathfrak m_R S$ are regular local rings,
\item $A = R/I$ and $B = S/J$ for some ideals $I$, $J$,
\item $J \subset S$ and
$\overline{J} = J/\mathfrak m_R \cap J \subset \overline{S}$
are generated by regular sequences, and
\item $A \to B$ and $R \to S$ are flat.
\end{enumerate}
Then $I$ is generated by a regular sequence.
\end{lemma}

\begin{proof}
Set $\overline{B} = B/\mathfrak m_RB = B/\mathfrak m_AB$ so that
$\overline{B} = \overline{S}/\overline{J}$.
Let $f_1, \ldots, f_{\overline{c}} \in J$ be elements such that
$\overline{f}_1, \ldots, \overline{f}_{\overline{c}} \in \overline{J}$
form a regular sequence generating $\overline{J}$.
Note that $\overline{c} = \dim(\overline{S}) - \dim(\overline{B})$,
see Lemma \ref{lemma-ci-well-defined}.
By Lemma \ref{lemma-grothendieck-regular-sequence}
the ring $S/(f_1, \ldots, f_{\overline{c}})$ is flat
over $R$. Hence $S/(f_1, \ldots, f_{\overline{c}}) + IS$ is flat over $A$.
The map $S/(f_1, \ldots, f_{\overline{c}}) + IS \to B$ is therefore a
surjection of finite $S/IS$-modules flat over $A$ which
is an isomorphism modulo $\mathfrak m_A$, and hence an
isomorphism by Lemma \ref{lemma-mod-injective}. In other words,
$J = (f_1, \ldots, f_{\overline{c}}) + IS$.

\medskip\noindent
By Lemma \ref{lemma-ci-well-defined} again the ideal $J$ is
generated by a regular sequence of $c = \dim(S) - \dim(B)$ elements. Hence
$J/\mathfrak m_SJ$ is a vector space of dimension $c$.
By the description of $J$ above there exist
$g_1, \ldots, g_{c - \overline{c}} \in I$ such that
$J$ is generated by
$f_1, \ldots, f_{\overline{c}}, g_1, \ldots, g_{c - \overline{c}}$
(use Nakayama's Lemma \ref{lemma-NAK}). Consider the ring
$A' = R/(g_1, \ldots, g_{c - \overline{c}})$ and the surjection
$A' \to A$. We see from the above that
$B = S/(f_1, \ldots, f_{\overline{c}}, g_1, \ldots, g_{c - \overline{c}})$
is flat over $A'$ (as $S/(f_1, \ldots, f_{\overline{c}})$ is flat
over $R$). Hence $A' \to B$ is injective (as it is faithfully flat,
see Lemma \ref{lemma-local-flat-ff}).
Since this map factors through $A$ we get $A' = A$.
Note that $\dim(B) = \dim(A) + \dim(\overline{B})$, and
$\dim(S) = \dim(R) + \dim(\overline{S})$, see
Lemma \ref{lemma-dimension-base-fibre-equals-total}.
Hence $c - \overline{c} = \dim(R) -\dim(A)$ by elementary algebra.
Thus $I = (g_1, \ldots, g_{c - \overline{c}})$ is generated
by a regular sequence according to Lemma \ref{lemma-ci-well-defined}.
\end{proof}






